\chapter*{第 5 章\quad 管路、孔口、管嘴的水力计算·孔珑题选}
\addcontentsline{toc}{chapter}{第 5 章\quad 管路、孔口、管嘴的水力计算·孔珑题选}
\markboth{第 5 章\quad 管路、孔口、管嘴的水力计算·孔珑题选}{}

\begin{tishi}
\textbf{题源：}本章题目取自孔珑主编《工程流体力学（第四版）》配套辅导书
《工程流体力学 学习指导及习题解答》（陈洁、袁铁江 编著，清华大学出版社）第 6 章
"管内流动和水力计算 液体出流"的课后习题解答段（原书书页 78--103，原书题号 6-1 至 6-41）。
孔珑该章把\textbf{管内流动阻力与孔口、管嘴出流合为一章}，与西华 818 考纲所用赵琴《流体力学与流体机械》
\textbf{第 5 章"管路、孔口、管嘴的水力计算"}一一对应，故本册原样归入第 5 章。

\begin{tabularx}{\linewidth}{@{}lX@{}}
\toprule
原书题号 & 处理方式 \\
\midrule
6-1 $\sim$ 6-29、6-32、6-33、6-35 $\sim$ 6-41 & \textbf{本章收录}（$\Rey$ 与流态判别、沿程与局部损失、串联与并联管路、管网、孔口与管嘴出流、非定常出流）\\
6-30 & \textbf{本章不收}（水泵特性曲线与管网工况点，属赵琴第 11--14 章流体机械，复试范围）\\
6-31 & \textbf{本章不收}（原书该题只有"略"字，既无题面也无解答）\\
6-34 & \textbf{本章不收}（离心泵吸水管最大安装高度，属泵的汽蚀与安装计算，复试范围）\\
\bottomrule
\end{tabularx}

\textbf{全章判据常量（必记）：}圆管下临界雷诺数 $\Rey_c=2000$（$\Rey<2000$ 层流，$\Rey>2000$ 紊流）；
圆管层流 $\lambda=64/\Rey$（与粗糙度无关）；紊流水力光滑区 $\lambda=0.3164/\Rey^{0.25}$（布拉修斯）；
圆管锐边进口 $\zeta=0.5$、出口流入大容器 $\zeta=1.0$；薄壁小孔口流量系数 $\mu\approx0.62$、
圆柱形外管嘴 $\mu\approx0.82$；管嘴正常工作的条件为 $H_0\le9\ \mathrm{m}$ 且 $l=(3\sim4)d$。
选择题中反复出现的"管嘴正常条件""虹吸管真空度""圆管流态判别"三类小计算，本册不再收编，见"教材精讲"分册第 5 章。

\textbf{做题要求：}① 凡求损失大小的题，一律\textbf{先算 $\Rey$ 判流态}，再选对应的 $\lambda$ 公式或查莫迪图；
② 伯努利类题必须写明\textbf{选断面、基准面、压强基准（相对/绝对）}；③ 串联管路\textbf{流量相等、损失相加}，
并联管路\textbf{损失相等、流量相加}；④ 全部题目都要有实质解答，不得写"略"。
原书题图一律用〔图示〕文字转述已知量，\textbf{不重绘图、不编造尺寸}；凡影印不清、题面数据只能由解答反推者，在章末按题号集中说明。
\end{tishi}

\begin{ti}\sloppy\emergencystretch=2em

\item[\textbf{6-1}] 半径为 $r_0$ 的圆管中流动为层流，流速恰好等于断面平均流速处与管轴之间的距离 $r$ 等于多少？
\xstarfull{5}\quad\yiju{E4 2015 年试题选择第 10 题（圆管层流过流断面切应力分布）；E4 2015 年试题填空第 5 题（层流断面平均流速＝最大流速的 0.5 倍）；E4 2022 年单选第 10 题（层流管径减半、压强损失 16 倍）；E1 slide33「第五章＝考试计算题重点」}\quad\nandu{易}

\item[\textbf{6-2}] 沿直径 $d=200\ \mathrm{mm}$ 的管道输送润滑油，质量流量 $q_m=9000\ \mathrm{kg/h}$，
油的密度 $\rho=900\ \mathrm{kg/m^3}$，运动黏度冬季为 $\nu_1=0.0001092\ \mathrm{m^2/s}$、
夏季为 $\nu_2=0.0000355\ \mathrm{m^2/s}$。试分别判断冬、夏两季润滑油在管中的流动状态。
\xstarfull{4}\quad\yiju{E4 2022 年单选第 9 题（两圆管流态判别）；E4 2015 年试题填空第 6 题（$\Rey$ 判流态）；E1 slide33「第五章＝考试计算题重点」}\quad\nandu{易}

\item[\textbf{6-3}] 蒸汽轮机的凝汽器有 400 条管径 $d=20\ \mathrm{mm}$ 的黄铜管，
管内循环流动着冷却水。为使散热迅速，管内需要形成稳定的紊流（$\Rey=33000$），
试求温度 10$^\circ$C 的冷却水的质量流量。
\tuyi{原书凝汽器示意图：管束由 400 条 $d=20\ \mathrm{mm}$ 的黄铜管组成，冷却水在管内流动、蒸汽在管外凝结放热。}
\xstarfull{4}\quad\yiju{E4 2022 年单选第 9 题（流态判别）；E4 2015 年试题填空第 6 题（$\Rey$ 判流态）；E1 slide33「第五章＝考试计算题重点」；E6 教材 \S 6.2 雷诺数与临界雷诺数}\quad\nandu{中}

\item[\textbf{6-4}] 欲使镀锌铁管内 $\Rey=3.5\times10^5$ 的流动是水力光滑的，管子的直径至少等于多少？
\xstarfull{3}\quad\yiju{E1 slide33「第五章＝考试计算题重点」；E6 教材 \S 6.4 尼古拉兹实验与水力光滑管判据；真题未直接考"光滑管所需最小管径"，故三星}\quad\nandu{中}

\item[\textbf{6-5}] 用水在直径 $d=30\ \mathrm{cm}$ 的水平管道上做沿程损失实验，
在相距 $l=120\ \mathrm{m}$ 的两点用水银压差计测量压强差，试求该管段的沿程阻力系数 $\lambda$。
（\textbf{原书该题题面影印不清、流量数据脱漏}，本册保留已知信息，解答按标准解法给出。）
\xstarfull{4}\quad\yiju{E4 2022 年单选第 10 题（沿程损失随管径的变化）；E5 题库 \texttt{10.pdf} 选择第 10 题（圆管与方管沿程损失比 0.886）；E1 slide33「第五章＝考试计算题重点」；E6 教材式（6.15）达西公式}\quad\nandu{中}

\item[\textbf{6-6}] 当 $\Rey=10^5$ 时，直径多大的镀锌铁管的沿程阻力系数与直径 $d_2=30\ \mathrm{cm}$ 的铸铁管相同？
\xstarfull{3}\quad\yiju{E1 slide33「第五章＝考试计算题重点」；E6 教材 \S 6.4 莫迪图（$\lambda$ 由 $\Rey$ 与相对粗糙度共同决定）；真题未直接考，故三星}\quad\nandu{难}

\item[\textbf{6-7}] 运动黏度 $\nu=4\times10^{-5}\ \mathrm{m^2/s}$ 的流体在直径 $d=1\ \mathrm{cm}$ 的管内以 $v=4\ \mathrm{m/s}$ 的速度流动，
求每米管长上的沿程损失。
\xstarfull{4}\quad\yiju{E4 2022 年单选第 10 题（层流沿程损失与管径的定量关系）；E4 2015 年试题填空第 5 题（层流平均流速＝最大流速 0.5 倍）；E1 slide33「第五章＝考试计算题重点」}\quad\nandu{易}

\item[\textbf{6-8}] 喷水泉的喷嘴为一截头圆锥体，长 $l=0.5\ \mathrm{m}$，两端直径 $d_1=40\ \mathrm{mm}$、$d_2=20\ \mathrm{mm}$，
竖直放置。若把计示压强 $p_{1e}=9.807\times10^4\ \mathrm{Pa}$ 的水引入喷嘴，
喷嘴内的能量损失 $h_w=1.6\ \mathrm{m}$（水柱），不计空气阻力，求喷出的体积流量和射流的上升高度。
\tuyi{原书喷水泉示意图：截头圆锥形喷嘴铅垂放置，下端进口直径 $d_1=40\ \mathrm{mm}$、计示压强 $p_{1e}=9.807\times10^4\ \mathrm{Pa}$，上端出口直径 $d_2=20\ \mathrm{mm}$，水自出口铅垂向上射入大气，喷嘴长 $l=0.5\ \mathrm{m}$。}
\xstarfull{4}\quad\yiju{E1 slide16「计算题考伯努利方程」；E1 slide33「第五章＝考试计算题重点」；E6 教材 \S 4.4 总流伯努利方程与连续性方程联用}\quad\nandu{中}

\item[\textbf{6-9}] 输油管的直径 $d=150\ \mathrm{mm}$、长 $l=5000\ \mathrm{m}$，出口端比进口端高 $h=10\ \mathrm{m}$；
输送油的质量流量 $q_m=15489\ \mathrm{kg/h}$，油的密度 $\rho=859.4\ \mathrm{kg/m^3}$，
进口端油压 $p_1=49\times10^4\ \mathrm{Pa}$，沿程阻力系数 $\lambda=0.03$，求出口端的油压 $p_e$。
\tuyi{原书输油管示意图：水平距离 $l=5000\ \mathrm{m}$ 的输油管，出口端高于进口端 $h=10\ \mathrm{m}$，进口端油压 $p_1=49\times10^4\ \mathrm{Pa}$。}
\xstarfull{4}\quad\yiju{E1 slide16「计算题考伯努利方程」；E1 slide33「第五章＝考试计算题重点」；E6 教材式（6.15）达西公式与总流伯努利方程联用}\quad\nandu{中}

\item[\textbf{6-10}] 水管直径 $d=250\ \mathrm{mm}$、长 $l=300\ \mathrm{m}$，绝对粗糙度 $e=0.25\ \mathrm{mm}$；
已知体积流量 $q_V=95\ \mathrm{L/s}$，水的运动黏度 $\nu=1\times10^{-6}\ \mathrm{m^2/s}$，求管段的沿程损失（米水柱）。
\xstarfull{5}\quad\yiju{E4 2022 年单选第 10 题（沿程损失计算）；E5 题库 \texttt{10.pdf} 选择第 10 题（沿程损失比 0.886）；E1 slide33「第五章＝考试计算题重点」；E2「第五章为主要计算考点」}\quad\nandu{中}

\item[\textbf{6-11}] 加热炉消耗 $q_m=300\ \mathrm{kg/h}$ 的重油，密度 $\rho=880\ \mathrm{kg/m^3}$，
运动黏度 $\nu=0.000025\ \mathrm{m^2/s}$。压力油箱位于喷油器轴线以上 $h=8\ \mathrm{m}$，
输油管直径 $d=25\ \mathrm{mm}$、长 $l=30\ \mathrm{m}$，求喷油器前重油的计示压强。
\xstarfull{4}\quad\yiju{E4 2022 年单选第 10 题（层流沿程损失）；E1 slide33「第五章＝考试计算题重点」；E6 教材 \S 6.3 层流 $\lambda=64/\Rey$ 与达西公式}\quad\nandu{中}

\item[\textbf{6-12}] 发动机润滑油的用量 $q_V=0.4\ \mathrm{cm^3/s}$，油从压力油箱经一输油管供给，
输油管直径 $d=6\ \mathrm{mm}$、长 $l=5\ \mathrm{m}$，油的密度 $\rho=820\ \mathrm{kg/m^3}$、
运动黏度 $\nu=0.000015\ \mathrm{m^2/s}$。设输油管终端压强等于大气压，求压力油箱所需的位置高度 $h$。
\xstarfull{4}\quad\yiju{E4 2022 年单选第 10 题（层流沿程损失）；E1 slide33「第五章＝考试计算题重点」；E6 教材 \S 6.3（层流）}\quad\nandu{中}

\item[\textbf{6-13}] 15$^\circ$C 的空气流过长 $l=200\ \mathrm{m}$、直径 $d=1.25\ \mathrm{m}$、绝对粗糙度 $e=1\ \mathrm{mm}$ 的管道，
已知沿程损失 $h_f=8\ \mathrm{cm}$（水柱），求空气的体积流量。
\xstarfull{3}\quad\yiju{E1 slide33「第五章＝考试计算题重点」；E6 教材 \S 6.3 达西公式（气流沿程损失按相当水柱换算后仍用达西公式）；真题未直接考气体管路，故三星}\quad\nandu{难}

\item[\textbf{6-14}] 用长 $l=15\ \mathrm{m}$、直径 $d=12\ \mathrm{mm}$ 的低碳钢管排出油箱中的油，
已知油的密度 $\rho=815.8\ \mathrm{kg/m^3}$、动力黏度 $\mu=0.01\ \mathrm{Pa\cdot s}$，
油面比管道出口高 $H=2\ \mathrm{m}$，求油的体积流量。
\xstarfull{4}\quad\yiju{E4 2022 年单选第 10 题（层流管流损失）；E1 slide33「第五章＝考试计算题重点」；E6 教材式（5.21）哈根--泊肃叶流量公式}\quad\nandu{中}

\item[\textbf{6-15}] 密度 $\rho=860\ \mathrm{kg/m^3}$、运动黏度 $\nu=5\times10^{-6}\ \mathrm{m^2/s}$ 的轻柴油，
通过长 $l=150\ \mathrm{m}$、绝对粗糙度 $e=0.45\ \mathrm{mm}$ 的铸铁管道从油池输送到储油库；
出油端比吸入端高 $H=25\ \mathrm{m}$，油泵能产生的压强 $p_1=3.43\times10^5\ \mathrm{Pa}$，
只计沿程损失，求必需的管道直径 $d$。（\textbf{原书题面流量数字影印不清}，本册按原书解答反推的体积流量计算。）
\xstarfull{4}\quad\yiju{E4 2022 年单选第 10 题（沿程损失与管径的定量关系）；E1 slide33「第五章＝考试计算题重点」；E6 教材 \S 6.3 达西公式与莫迪图迭代}\quad\nandu{难}

\item[\textbf{6-16}] 用新铸铁管输送 25$^\circ$C 的水，体积流量 $q_V=300\ \mathrm{L/s}$，
在长 $l=1000\ \mathrm{m}$ 的管道上沿程损失为 $h_f=2\ \mathrm{m}$（水柱），求必需的管道直径。
\xstarfull{4}\quad\yiju{E4 2022 年单选第 10 题（沿程损失与管径的定量关系）；E1 slide33「第五章＝考试计算题重点」；E6 教材 \S 6.3 达西公式与莫迪图迭代}\quad\nandu{难}

\item[\textbf{6-17}] 已知油的密度 $\rho=800\ \mathrm{kg/m^3}$、动力黏度 $\mu=0.069\ \mathrm{Pa\cdot s}$，
在连接两容器的光滑管中流动（图 6-56），$H=3\ \mathrm{m}$；管长 $l=30\ \mathrm{m}$、直径 $d=30\ \mathrm{cm}$。
计及沿程和局部损失时，管内的流量为多少？
\tuyi{原书图 6-56：两个敞口容器，液面高差 $H=3\ \mathrm{m}$，其间用一条长 $l=30\ \mathrm{m}$、直径 $d=30\ \mathrm{cm}$ 的光滑管连接，油自高液面容器流入低液面容器。}
\xstarfull{4}\quad\yiju{E4 2015 年试题计算第 3 题（水头损失类计算大题）；E1 slide33「第五章＝考试计算题重点」；E6 教材 \S 6.5 局部损失与 \S 6.3 布拉修斯公式迭代}\quad\nandu{难}

\item[\textbf{6-18}] 在习题 6-17 的管道中安装一个阀门，使管内流量减少到原流量的一半，
求阀门的局部损失系数以及按该管道换算的等值长度。
\xstarfull{4}\quad\yiju{E4 2015 年试题计算第 3 题（局部水头损失的计算）；E1 slide33「第五章＝考试计算题重点」；E6 教材 \S 6.5 局部损失系数与等值长度}\quad\nandu{中}

\item[\textbf{6-19}] 水温对应的运动黏度 $\nu=1.51\times10^{-6}\ \mathrm{m^2/s}$，体积流量 $q_V=15\ \mathrm{m^3/h}$ 的水，
在直径 $d=50\ \mathrm{mm}$、绝对粗糙度 $e=0.2\ \mathrm{mm}$ 的 90$^\circ$ 弯管中流动；
水银压差计两接点相距 $l=0.8\ \mathrm{m}$，读数 $h=20\ \mathrm{mm}$，求弯管的局部损失系数。
\xstarfull{4}\quad\yiju{E4 2015 年试题计算第 3 题「突然缩小的水头损失」（局部水头损失的实测与计算）；E1 slide33「第五章＝考试计算题重点」；E6 教材 \S 6.5 局部损失与差压计测损失}\quad\nandu{中}

\item[\textbf{6-20}] 管径由 $d_1=50\ \mathrm{mm}$ 突然扩大到 $d_2=100\ \mathrm{mm}$，流过体积流量 $q_V=16\ \mathrm{m^3/h}$ 的水；
在截面改变处插入内充四氯化碳（$\rho'=1600\ \mathrm{kg/m^3}$）的压差计，读得液面高差 $h=173\ \mathrm{mm}$。
求管径扩大处的局部损失系数，并把结果与理论值相比较。
\xstarfull{5}\quad\yiju{E4 2015 年试题计算第 3 题「突然缩小的水头损失」（与突扩同为包达--卡诺局部损失，同一考点）；E1 slide33「第五章＝考试计算题重点」；E2「第五章为主要计算考点」；E6 教材 \S 6.5 突然扩大局部损失式（6.35）}\quad\nandu{中}

\item[\textbf{6-21}] 鼓风机站供给高炉车间的空气量 $q_V=120000\ \mathrm{m^3/h}$，空气温度 $t=20^\circ\mathrm{C}$、
运动黏度 $\nu=0.0000157\ \mathrm{m^2/s}$；输气管总长 $l=120\ \mathrm{m}$、绝对粗糙度 $e=0.5\ \mathrm{mm}$，
管上有 5 个弯曲半径 $R_1=2.6\ \mathrm{m}$ 的弯曲处、4 个弯曲半径 $R_2=1.3\ \mathrm{m}$ 的弯曲处和 2 个 $\zeta=2.5$ 的闸阀；
管内空气流速 $u=25\ \mathrm{m/s}$，热风炉进口处的计示压强 $p_{ie}=156896\ \mathrm{Pa}$。
求输气管的管径与鼓风机出口处的计示压强。
\tuyi{原书输气管示意图：鼓风机出口经总长 $l=120\ \mathrm{m}$ 的输气管送往热风炉，管上有 5 个 $R_1=2.6\ \mathrm{m}$ 的弯曲处、4 个 $R_2=1.3\ \mathrm{m}$ 的弯曲处和 2 个闸阀（$\zeta=2.5$）；管内空气流速 $u=25\ \mathrm{m/s}$，热风炉进口计示压强 $p_{ie}=156896\ \mathrm{Pa}$。}
\xstarfull{3}\quad\yiju{E1 slide33「第五章＝考试计算题重点」；E6 教材 \S 6.5 弯管与阀门局部损失系数、\S 6.4 粗糙管平方阻力区；本题风机只作压强源，未涉及风机特性曲线，真题未直接考，故三星}\quad\nandu{难}

\item[\textbf{6-22}] 在三路管状空气预热器中，将质量流量 $q_m=5816\ \mathrm{kg/h}$ 的空气从 $t_1=20^\circ\mathrm{C}$ 加热到 $t_2=160^\circ\mathrm{C}$；
预热器高 $H=4\ \mathrm{m}$，管系损失系数 $\zeta=6$（对管内平均流速而言），管系截面积 $A_1=0.4\ \mathrm{m^2}$，
连接箱截面积 $A_2=0.8\ \mathrm{m^2}$，管径与拐弯处曲率半径之比 $d/R=1$；不计沿程损失，
试按空气的平均温度求流经空气预热器的总压降。
\xstarfull{3}\quad\yiju{E1 slide33「第五章＝考试计算题重点」；E6 教材 \S 6.5 局部损失系数（对平均流速）；本题为局部损失与浮升压降的简单叠加，真题未直接考，故三星}\quad\nandu{难}

\item[\textbf{6-23}] 为测定新阀门的压强降，在一钢管管系上装设水银压差计。
钢管内径均为 $50\ \mathrm{mm}$、绝对粗糙度 $e=0.44\ \mathrm{mm}$，弯管 $d/R=0.1$；
水泵流量稳定在 $q_V=12\ \mathrm{m^3/h}$，水银压差计读数 $150\ \mathrm{mm}$，
上游容器水面高出管道出口 $1.8\ \mathrm{m}$、出口低于水面 $2\ \mathrm{m}$。试求：
\begin{xiaowen}
\item 阀门的压强降 $\Delta p$；
\item 阀门的局部损失系数 $\zeta$；
\item 阀门前的计示压强；
\item 水泵供给水的功率。
\end{xiaowen}
\xstarfull{3}\quad\yiju{E4 2015 年试题计算第 3 题（局部损失实测与系数换算）；E5 题库 \texttt{10.pdf} 计算第 2 题（水箱与垂直管道出流，求 $Q$ 与管长的关系）；E1 slide33「第五章＝考试计算题重点」}\quad\nandu{难}

\item[\textbf{6-24}] 中压锅炉的过热器由长 $L=24\ \mathrm{m}$、直径 $d=35\ \mathrm{mm}$、绝对粗糙度 $e=0.08\ \mathrm{mm}$ 的无缝钢管组成，
管上有 13 个弯头（每个 $\zeta=0.2$）、进口 $\zeta=0.7$、出口 $\zeta=1.1$；蒸汽的体积流量 $q_V=0.026\ \mathrm{m^3/s}$、
温度 $423^\circ\mathrm{C}$、运动黏度 $\nu=1.98\times10^{-6}\ \mathrm{m^2/s}$、密度 $\rho=12.8\ \mathrm{kg/m^3}$，
出口压强 $p_2=3.923\times10^6\ \mathrm{Pa}$，求过热器进口处的压强 $p_1$。
\xstarfull{4}\quad\yiju{E4 2015 年试题计算第 3 题（沿程与局部损失叠加求压强）；E1 slide33「第五章＝考试计算题重点」；E6 教材 \S 6.3 粗糙管平方阻力区与 \S 6.5 局部损失}\quad\nandu{中}

\item[\textbf{6-25}] 车间输水管道由两段串联组成，绝对粗糙度 $e=0.5\ \mathrm{mm}$；
第一段 $l_1=800\ \mathrm{m}$、$d_1=150\ \mathrm{mm}$，第二段 $l_2=600\ \mathrm{m}$、$d_2=125\ \mathrm{mm}$；
压力水塔水头 $H=20\ \mathrm{m}$，不计局部损失，求阀门全开时通过的最大流量。
\xstarfull{4}\quad\yiju{E1 slide33「第五章＝考试计算题重点」；E2「第五章同为主要计算考点」；E4 2022 年单选第 10 题（沿程损失）；E6 教材 \S 6.6 串联管路}\quad\nandu{中}

\item[\textbf{6-26}] 两容器用两段新低碳钢管串联连接（图 6-62）：$l_1=30\ \mathrm{m}$、$d_1=20\ \mathrm{cm}$，
$l_2=60\ \mathrm{m}$、$d_2=30\ \mathrm{cm}$；第一段为锐边进口，第二段上装有阀门（$\zeta=3.5$）；
体积流量 $q_V=0.2\ \mathrm{m^3/s}$，求必需的总水头 $H$。
\tuyi{原书图 6-62：两敞口容器间用两段钢管串联，第一段 $l_1=30\ \mathrm{m}$、$d_1=20\ \mathrm{cm}$ 自容器 1 锐边进口；第二段 $l_2=60\ \mathrm{m}$、$d_2=30\ \mathrm{cm}$ 上装阀门，出口流入容器 2。}
\xstarfull{4}\quad\yiju{E4 2015 年试题计算第 3 题（突扩／突缩局部损失）；E1 slide33「第五章＝考试计算题重点」；E6 教材 \S 6.6 串联管路（流量相等、损失相加）}\quad\nandu{中}

\item[\textbf{6-27}] 两容器间用两根并联铸铁管连接：$l_1=2500\ \mathrm{m}$、$d_1=1.2\ \mathrm{m}$；
$l_2=2000\ \mathrm{m}$、$d_2=1\ \mathrm{m}$；两管绝对粗糙度均为 $e=0.00045\ \mathrm{m}$。
总水头 $H=3.6\ \mathrm{m}$，求 20$^\circ$C 水的总流量。
\tuyi{原书并联管路示意图：两容器（上、下游水头差 $H=3.6\ \mathrm{m}$）之间并联两根铸铁管，管 1 长 $2500\ \mathrm{m}$、直径 $1.2\ \mathrm{m}$，管 2 长 $2000\ \mathrm{m}$、直径 $1\ \mathrm{m}$，两管进口、出口分别位于同一断面。}
\xstarfull{4}\quad\yiju{E1 slide33「第五章＝考试计算题重点」；E2「第五章同为主要计算考点」；E6 教材 \S 6.6 并联管路（损失相等、流量相加）}\quad\nandu{中}

\item[\textbf{6-28}] 总流量 $q_V=25\ \mathrm{L/s}$ 的输水管接入两条并联管道：$l_1=500\ \mathrm{m}$、$d_1=10\ \mathrm{cm}$、$e_1=0.2\ \mathrm{mm}$；
$l_2=900\ \mathrm{m}$、$d_2=15\ \mathrm{cm}$、$e_2=0.5\ \mathrm{mm}$。求两支管的流量分配和并联段进出口的水头损失。
\xstarfull{4}\quad\yiju{E1 slide33「第五章＝考试计算题重点」；E2「第五章同为主要计算考点」；E6 教材 \S 6.6 并联管路流量分配（试算迭代）}\quad\nandu{难}

\item[\textbf{6-29}] 物料烘干器的并联蛇形管：粗管 $l_1=8\ \mathrm{m}$、$d_1=25\ \mathrm{mm}$、$e_1=0.25\ \mathrm{mm}$、$\zeta_1=4.2$；
细管 $l_2=14\ \mathrm{m}$、$d_2=10\ \mathrm{mm}$、$e_2=0.11\ \mathrm{mm}$、$\zeta_2=5.2$。
蒸汽质量流量 $q_m=0.01\ \mathrm{kg/s}$、温度 $180^\circ\mathrm{C}$、密度 $\rho=5\ \mathrm{kg/m^3}$、
动力黏度 $\mu=15.1\times10^{-6}\ \mathrm{Pa\cdot s}$，不计压缩性，求两管的流量与管系的压降。
\xstarfull{3}\quad\yiju{E1 slide33「第五章＝考试计算题重点」；E6 教材 \S 6.6 并联管路 与 \S 6.5 局部损失；本题为蒸汽（密度小）并联管路的迭代，真题未直接考，故三星}\quad\nandu{难}

\item[\textbf{6-32}] 图 6-66 所示管网中 $AB=AC=150\ \mathrm{m}$、$BD=CE=250\ \mathrm{m}$、$AD=AE=DF=EF=300\ \mathrm{m}$，
各管直径均为 $0.4\ \mathrm{m}$，$F$ 点的计示压强为 $50\ \mathrm{kPa}$，各管沿程阻力系数均为 $\lambda=0.03$，
忽略局部损失。管网各节点的分出流量见图（$A$ 点流入 $0.4\ \mathrm{m^3/s}$，$B$、$C$ 各流入 $0.2\ \mathrm{m^3/s}$，
$D$、$E$ 各流出 $0.3\ \mathrm{m^3/s}$，$F$ 流出 $0.2\ \mathrm{m^3/s}$）。求各管段的流量与 $A$、$B$、$C$ 三点的压强。
\tuyi{原书图 6-66：环状管网，环路由 $A$、$B$、$D$、$F$、$E$、$C$ 构成，$AB=AC=150\ \mathrm{m}$、$BD=CE=250\ \mathrm{m}$、$AD=AE=DF=EF=300\ \mathrm{m}$，各管直径均为 $0.4\ \mathrm{m}$；$F$ 点为计示压强 $50\ \mathrm{kPa}$ 的节点。}
\xstarfull{3}\quad\yiju{E1 slide33「第五章＝考试计算题重点」；E6 教材 \S 6.6 管网计算（节点连续性方程与回路方程联立）；真题未直接考环状管网，故三星}\quad\nandu{难}

\item[\textbf{6-33}] 用虹吸管将油从容器中引出。虹吸管由两根直径 $d=10\ \mathrm{cm}$ 的光滑管用靠得很近的回转管连接，
总长 $l=11.4\ \mathrm{m}$；末端 $A$ 接一段长 $15\ \mathrm{cm}$、出口直径 $d_2=5\ \mathrm{cm}$ 的喷嘴，
其损失按 $0.1v^2/(2g)$ 计；油的密度 $\rho=800\ \mathrm{kg/m^3}$、动力黏度 $\mu=0.01\ \mathrm{Pa\cdot s}$。
求虹吸管的流量以及 $A$、$B$ 两点的计示压强。
\tuyi{原书虹吸管示意图：\ 油自敞口容器经虹吸管翻越而出，最高点记为 $B$（高于容器液面）；管末端 $A$ 接一段长 $15\ \mathrm{cm}$、出口直径 $d_2=5\ \mathrm{cm}$ 的喷嘴，油自喷嘴排入大气；虹吸管由两根 $d=10\ \mathrm{cm}$ 的光滑管经靠得很近的回转管连接，管段长分别为 $6$、$1.8$、$3.6\ \mathrm{m}$，合计 $l=11.4\ \mathrm{m}$。}
\xstarfull{5}\quad\yiju{E4 2022 年单选第 7 题（虹吸管真空度）；E5 题库选择第 8 题（虹吸管最高处压强小于大气压）；E1 slide33「第五章＝考试计算题重点」；E6 教材 \S 6.5 虹吸管计算}\quad\nandu{难}

\item[\textbf{6-35}] 封闭大水箱经直径 $d=12.5\ \mathrm{mm}$ 的薄壁小孔口定常出流，孔口中心以上的水头 $H=1.8\ \mathrm{m}$，
体积流量 $q_V=1.5\ \mathrm{L/s}$，流量系数 $C_d=0.6$，求液面上气体的计示压强。
\xstarfull{4}\quad\yiju{E4 2026 回忆版计算第 5 题（孔口／管嘴的体积流量，同属书上例题）；E1 slide33「第五章＝考试计算题重点」；E6 教材 \S 6.7 孔口出流基本公式}\quad\nandu{中}

\item[\textbf{6-36}] 薄壁容器上有一个锐缘小孔口，直径 $d=7.5\ \mathrm{mm}$，孔口中心线以上的水头 $H=1.2\ \mathrm{m}$，
测得孔口出流的质量流量为 $8\ \mathrm{kg/min}$；又测得射流轨迹上某点距孔口水平距离 $x=1.25\ \mathrm{m}$、竖直下落距离 $z=0.35\ \mathrm{m}$。
求孔口的流速系数、流量系数、收缩系数和损失系数。
\xstarfull{5}\quad\yiju{E4 2026 回忆版计算第 5 题（管嘴／孔口体积流量，书上例题同型）；E4 2022 年单选第 8 题（管嘴正常工作的流量系数口径）；E1 slide33「第五章＝考试计算题重点」}\quad\nandu{中}

\item[\textbf{6-37}] 密度 $\rho=900\ \mathrm{kg/m^3}$ 的油从直径 $d=2\ \mathrm{cm}$ 的孔口射出，
射流冲击到孔口外的挡板上，冲击力 $F=20\ \mathrm{N}$；孔口前油液的计示压强为 $45000\ \mathrm{Pa}$，
出口的体积流量为 $2.29\ \mathrm{L/s}$。求孔口的流速系数、流量系数和收缩系数。
\xstarfull{4}\quad\yiju{E1 slide16「计算题考动量方程」；E1 slide33「第五章＝考试计算题重点」；E6 教材 \S 6.7 孔口出流系数与 \S 4.5 动量方程}\quad\nandu{中}

\item[\textbf{6-38}] 水箱侧壁装一缩放管嘴，喉部直径 $d_1=4\ \mathrm{cm}$，水头 $H=3\ \mathrm{m}$，
大气压 $p_a=10.33\ \mathrm{m}$（水柱），喉部压强为 $2.5\ \mathrm{m}$（水柱）；收缩段损失忽略，
渐扩段的损失按从直径 $d_1$ 突扩至 $d_2$ 的损失的 $20\%$ 计。
求喉部的流速与流量、出口处的流速以及出口截面直径 $d_2$。
\tuyi{原书缩放管嘴示意图：水箱侧壁装一缩放管嘴，收缩段末端为喉部（直径 $d_1=4\ \mathrm{cm}$、压强 $2.5\ \mathrm{m}$ 水柱），其后渐扩段出口直径 $d_2$ 待求；水箱内孔口中心以上的水头 $H=3\ \mathrm{m}$，大气压 $p_a=10.33\ \mathrm{m}$ 水柱。}
\xstarfull{5}\quad\yiju{E4 2022 年单选第 8 题（圆柱形直角外管嘴的正常工作条件）；E4 2026 回忆版计算第 5 题（管嘴体积流量，书上例题）；E1 slide33「第五章＝考试计算题重点」}\quad\nandu{难}

\item[\textbf{6-39}] 水箱中恒定水深 $H=5\ \mathrm{m}$，铅直管直径 $d=20\ \mathrm{cm}$。已知大气压 $p_a=1.013\times10^5\ \mathrm{Pa}$，
水的饱和压强 20$^\circ$C 时为 $2340\ \mathrm{Pa}$、60$^\circ$C 时为 $20000\ \mathrm{Pa}$。
忽略损失，为使铅直管进口处不出现气穴，求水温分别为 20$^\circ$C 和 60$^\circ$C 时的最大允许管长和最大理论流量。
\tuyi{原书图 6-70：水箱中恒定水深 $H=5\ \mathrm{m}$，箱底接一根直径 $d=20\ \mathrm{cm}$ 的铅直出流管，重力作用下自由出流至大气；要求管进口处不出现气穴（进口压强不低于该温度下的饱和压强）。}
\xstarfull{3}\quad\yiju{E4 2022 年单选第 7 题（虹吸管真空度限制，与本题"进口压强不得低于饱和压强"同属真空度类判据）；E1 slide33「第五章＝考试计算题重点」；E6 教材 \S 6.1 气穴与气蚀}\quad\nandu{难}

\item[\textbf{6-40}] 装满水的锥台形容器（图 6-71），$d_1=1\ \mathrm{m}$、$d_2=2\ \mathrm{m}$、$H=2\ \mathrm{m}$，顶盖通大气，
底部中心有一直径 $d=2.5\ \mathrm{cm}$ 的小孔口，流量系数 $C_d=0.60$。求通过小孔口将水放空所需要的时间。
\tuyi{原书图 6-71：锥台形容器，顶盖通大气，上口直径 $d_1=1\ \mathrm{m}$、底口直径 $d_2=2\ \mathrm{m}$、高（水深）$H=2\ \mathrm{m}$，底部中心开 $d=2.5\ \mathrm{cm}$ 的小孔口，水自孔口泄空。}
\xstarfull{3}\quad\yiju{E1 slide33「第五章＝考试计算题重点」；E6 教材 \S 6.7 变水头下的孔口出流与放空时间积分；真题未直接考放空时间，故三星}\quad\nandu{中}

\item[\textbf{6-41}] 图 6-72 所示两连通容器，液体从左侧容器流向右侧容器，
求两容器液位差从 $H_1$ 变到 $H_2$ 所需要的时间。
\tuyi{原书图 6-72：甲、乙两容器（截面面积分别为 $A_1$、$A_2$）底部用短管（面积 $a$、流量系数 $C_d$）连通，液体自左侧容器流向右方容器；初始液位差为 $H_1$，终了液位差为 $H_2$。}
\xstarfull{2}\quad\yiju{E1 slide33「第五章＝考试计算题重点」；E6 教材 \S 6.7 变水头孔口出流（连通容器液位差随时间变化）；真题未出现这一形式，属孔口非定常出流的延伸，故二星}\quad\nandu{难}

\end{ti}

\begin{tishi}
\textbf{〔收录与不收录的对照〕}孔珑该章共 41 题（6-1 至 6-41），本章收录 \textbf{38 题}：
6-1 $\sim$ 6-29（共 29 题）、6-32、6-33、6-35 $\sim$ 6-41（共 9 题）。
\textbf{不收 3 题}：6-30（原书图 6-64 给出水泵特性表，要求解水泵工况点与通过流量——属泵与管路联合运行，
归赵琴第 11--14 章流体机械，复试范围）；6-31（原书只有"略"一字）；6-34（离心式水泵吸水管最大允许安装高度，
需用泵进口的汽蚀余量与饱和压强，属泵的汽蚀计算，复试范围）。

\textbf{〔原书各题对应的书页〕}6-1 与 6-2 起于 p78；6-3、6-4、6-5 在 p79；6-6 $\sim$ 6-9 在 p80；
6-10、6-11 在 p81；6-12 $\sim$ 6-14 在 p82；6-15 跨 p83--84；6-16、6-17 在 p84--85；6-18、6-19 在 p85--86；
6-20、6-21 在 p87--88；6-22、6-23 在 p88--90；6-24 在 p90；6-25 在 p91；6-26、6-27 在 p91--92；
6-28、6-29 在 p92--95；6-32、6-33 在 p96--98；6-35 $\sim$ 6-37 在 p99--100；
6-38 跨 p100--101；6-39 跨 p101--102；6-40、6-41 在 p102--103。

\textbf{〔影印不清处（解题时请以题面与题图为准）〕}
\textbf{第 6-5 题：}原书该页题面与解答均影印不清，只能读出"用水在 $d=30\ \mathrm{cm}$ 水平管道上做沿程损失实验、
相距 $120\ \mathrm{m}$ 两点用水银压差计测量"与解答中的 $h_f=\rho_{Hg}gh$（$h=0.33\ \mathrm{m}$）、
$A=\pi\times0.15^2$、"$=203.53$"等碎片，\textbf{实验流量数据完全脱漏}。本册按标准解法给出方法与封闭式，
并把原书残留数据列在解答中，请以你书上的题面数据代入。\\
\textbf{第 6-6 题：}原书解答中"由莫迪图可知"以后的数值被切断，本册按"$\Rey$ 与相对粗糙度同时相等则 $\lambda$ 相等"复原。\\
\textbf{第 6-15 题：}原书题面流量印作"$q_m=10\ \mathrm{kg/h}$"，与原书解答反推的体积流量 $q_V=0.0323\ \mathrm{m^3/s}$ 不符，
本册按解答口径（约 $100\ \mathrm{t/h}$）计算，见章末存疑。\\
\textbf{第 6-23 题：}原书该题题面为一段长文字描述加多条件，OCR 后语序混乱；本册按原书解答反推的条件
（钢管 $d=50\ \mathrm{mm}$、$e=0.44\ \mathrm{mm}$、弯管 $d/R=0.1$、$q_V=12\ \mathrm{m^3/h}$、水银压差计 $150\ \mathrm{mm}$、
水位高出出口 $1.8\ \mathrm{m}$、出口低于水面 $2\ \mathrm{m}$）转述，若与你的题面略有出入，方法与解答完全一致。\\
\textbf{第 6-32 题：}管网各节点的分出流量由原书解答用连续性方程反推（$A$ 流入 $0.4$、$B$ 与 $C$ 各流入 $0.2$、
$D$ 与 $E$ 各流出 $0.3$、$F$ 流出 $0.2\ \mathrm{m^3/s}$），已在题面注明。\\
\textbf{第 6-41 题：}原书解答在 p103 末行被影印切断（推理写到 $\sqrt{H_1}-\sqrt{H_2}$ 处中断），
原书最终答案未印出，本册按标准解法补全并加提示框说明。
\end{tishi}

\begin{tishi}
\textbf{〔OCR 校订清单（量纲与指数缺字，逐条订正）〕}
\begin{xiaowen}
\item 第 6-12 题（p82）：润滑油的用量印作 $q_V=0.4\ \mathrm{cm^2/s}$，\textbf{应为 $0.4\ \mathrm{cm^3/s}$}（面积误作体积）。
\item 第 6-13 题（p82）、第 6-39 题（p101--102）：体积流量印作 $\mathrm{m^2/s}$，\textbf{一律应为 $\mathrm{m^3/s}$}。
\item 第 6-39 题（p101--102）：水的密度印作 $\mathrm{kg/m^2}$，\textbf{应为 $\mathrm{kg/m^3}$}。
\item 第 6-15 题（p83--84）：能量方程迭代式印作 $\lambda=1211.56\,d$，\textbf{应为 $\lambda=1211.56\,d^5$}（指数缺字，
由 $d=0.11\ \mathrm{m}$ 时 $\lambda=0.01951$ 可验算）。
\item 第 6-5 题（p79）：解答中 $h_f=\rho_{Hg}gh=43982.3$ 后误注单位为 $\mathrm{m}$，\textbf{该值为压强 $\mathrm{Pa}$}
（水头应为 $(\rho_{Hg}/\rho-1)h=4.488\ \mathrm{m}$）。
\item 第 6-35 题（p99）：液面气体的计示压强印作 $19007633\ \mathrm{Pa}$，小数点被吞掉，按流量 $1.5\ \mathrm{L/s}$ 反算
\textbf{应为 $1.9008\times10^5\ \mathrm{Pa}$}（即 $190076.33\ \mathrm{Pa}$，量级订正）。
\item 第 6-19 题（p86）：题面运动黏度印作 $0.000015\ \mathrm{m^2/s}$，按解答 $\Rey=70298$ 反推
\textbf{应为 $1.51\times10^{-6}\ \mathrm{m^2/s}$}（《水力学》表 15$^\circ$C 水的 $\nu=1.139\times10^{-6}$，
$5^\circ$C 为 $1.519\times10^{-6}$，本题按 1.51 计）。
\item 第 6-17 题（p84--85）、第 6-18 题（p85--86）：$\mu=0.69\ \mathrm{Pa\cdot s}$ 与 $\mu=0.069\ \mathrm{Pa\cdot s}$ 两处印刷不一，
按 $\Rey=19759$、$v=5.68\ \mathrm{m/s}$ 反推\textbf{取 $\mu=0.069\ \mathrm{Pa\cdot s}$}。
\item 第 6-36 题（p99）：解答页 OCR 乱码，四个系数按定义式复原（流速系数用射流轨迹公式、流量系数用质量流量公式）。
\item 第 6-3 题（p79）：10$^\circ$C 水的 $\rho=999.73\ \mathrm{kg/m^3}$、$\nu=1.308\times10^{-6}\ \mathrm{m^2/s}$ 为查表值，已补入。
\end{xiaowen}
\end{tishi}

\begin{tishi}
\textbf{〔存疑处（本册的处理口径）〕}
\begin{xiaowen}
\item 第 6-3 题：原书按 $400/2=200$ 条管子并联给出 $q_m=135.6\ \mathrm{kg/s}$（相当于双流程凝汽器每次只有 200 条管通水）；
若按题面"400 条管全部并联"计，应为 $271.1\ \mathrm{kg/s}$。本册两个口径并列，以原书答案为准。
\item 第 6-6 题：镀锌铁管的绝对粗糙度原书取 $e_1=0.39\ \mathrm{mm}$，与第 6-4 题取 $e=0.25\ \mathrm{mm}$ 不一致；
本册按原书各自取值，并按铸铁管 $e_2=0.25\sim0.42\ \mathrm{mm}$ 给出直径区间。
\item 第 6-15 题：题面流量（$10\ \mathrm{kg/h}$）与原书解答反推的 $q_V=0.0323\ \mathrm{m^3/s}$ 相差约三个量级，
本册按解答口径计算，结果为 $d\approx0.12\ \mathrm{m}$。
\item 第 6-22 题：原书连接箱损失的折算系数 $6\times0.291$ 与末式所用的 $h_w=6.1895\ \mathrm{m}$ 前后不一致，
本册按原书末式取 $\Delta p\approx97\ \mathrm{Pa}$，并在解答中给出两种口径。
\item 第 6-29、6-33 题：原书解答内部数值前后不一致（第 6-33 题同时印出 $\lambda=0.03164/\Rey^{0.25}$ 与 $\lambda=0.00289$），
本册按迭代收敛值取用，收敛过程写在解答里。
\item 第 6-39 题：本题只涉及铅直管的进口气穴判据（与泵无关），故按考纲收编；
若你的参考书上把气穴一节归入流体机械，可按"真空度限制"一类题处理。
\end{xiaowen}
\end{tishi}
